\originalsection{18.04第一次模拟考试}
\sourceinfo{Jeremy Orloff, Spring 2018. 题面 mit18\_04s18\_pex1.pdf, PDF pp.1--3; 官方解答 pex1-qa, pp.1--9. MIT OCW, CC BY-NC-SA 4.0. 原卷12道大题; 圆周及闭曲线未另指方向时沿官方解答取正向.}
\begin{exercise}\label{ass:pex1-1}
\textbf{原题 Pex1.1。} (a) 求 $(z+2)/(z-1)$ 的实部和虚部. (b) 解 $z^4-\ii=0$. (c) 求 $\sqrt{\sqrt{\ii}}$ 的所有可能值. (d) 用 $\cos x,\sin x$ 表示 $\cos4x$. (e) 三角不等式 $|z_1+z_2|\le|z_1|+|z_2|$ 何时取等? (f) 画图表示 $z$ 与 $1/z$ 的极坐标.
\end{exercise}
\begin{solution}
据官方解答翻译并补全零值条件. (a) 写 $z=x+\ii y\ne1$, 乘分子、分母以 $x-1-\ii y$, 得实部 $((x+2)(x-1)+y^2)/((x-1)^2+y^2)$, 虚部 $-3y/((x-1)^2+y^2)$.

(b) $\ii=\ee^{\ii(\pi/2+2k\pi)}$, 四根为 $\ee^{\ii(\pi/8+k\pi/2)}$, $k=0,1,2,3$. (c) 两次平方根各有两个选择, 得到的四数恰好满足 $z^4=\ii$, 故与(b)相同. 若只取两次主平方根则只有 $\ee^{\ii\pi/8}$, 不符合题目“所有可能值”.

(d) 取 $(\cos x+\ii\sin x)^4$ 的实部, 得 $\cos4x=\cos^4x-6\cos^2x\sin^2x+\sin^4x$. (e) 至少一数为零时取等; 两数非零时取等当且仅当 $z_1\bar z_2$ 是正实数, 即两数在同一条由原点发出的射线上. (f) 若 $z=r\ee^{\ii\theta}$, $r>0$, 则 $1/z=r^{-1}\ee^{-\ii\theta}$: 模互为倒数, 辐角相反. 图只用一组示例坐标表示这一关系.
\begin{center}\begin{tikzpicture}[scale=.9,>=Stealth]
\draw[->](-.4,0)--(3.6,0) node[right]{$\Re$};\draw[->](0,-1.6)--(0,2.4) node[above]{$\Im$};
\draw[thick,->](0,0)--(2.6,1.5) node[above]{$r\ee^{\ii\theta}$};
\draw[thick,->](0,0)--(.87,-.5) node[right]{$r^{-1}\ee^{-\ii\theta}$};
\draw(.65,0) arc(0:30:.65);\node at (.94,.23){$\theta$};\draw(.5,0) arc(0:-30:.5);\node at (.35,-.72){$-\theta$};
\end{tikzpicture}\end{center}
\end{solution}
\begin{exercise}\label{ass:pex1-2}
\textbf{原题 Pex1.2。} (a) 证明 $\sinh z=-\ii\sin(\ii z)$. (b) 用实三角函数及双曲函数求 $\cos z$ 的实部、虚部, 其中 $z=x+\ii y$. (c) 是否总有 $|a^b|=|a|^{|b|}$?
\end{exercise}
\begin{solution}
据官方解答翻译并补全. (a) $-\ii\sin(\ii z)=-(\ee^{-z}-\ee^z)/2=\sinh z$. (b) 由指数定义 $\cos(x+\ii y)=\cos x\cosh y-\ii\sin x\sinh y$, 故两部分分别是 $\cos x\cosh y$ 与 $-\sin x\sinh y$. (c) 错. 即使底数为正并以实对数定义 $a^b=\ee^{b\log a}$, 取 $a=\ee,b=\ii$, 左侧 $|\ee^\ii|=1$, 右侧 $\ee^{|\ii|}=\ee$.
\end{solution}
\begin{exercise}\label{ass:pex1-3}
\textbf{原题 Pex1.3。} (a) 证明 $f(z)=(z-\ii)/(z+\ii)$ 把上半平面映到单位圆盘: (i) 实轴映到单位圆周; (ii) $\ii$ 映到0; (iii) 据此说明上半平面映到单位圆盘. (b) 证明 $f(z)=(z+2)/(z-1)$ 把单位圆周映到直线 $\Re w=-1/2$.
\end{exercise}
\begin{solution}
据官方解答翻译, 单射与满射的验证为编者补充（AI）. (a)(i) 实数 $t$ 满足 $|t-\ii|=|t+\ii|$, 故 $|f(t)|=1$; 扩充平面中的无穷远点映到1. (ii) 直接代入得 $f(\ii)=0$. (iii) 对 $z=x+\ii y$, $y>0$, 有 $|z+\ii|^2-|z-\ii|^2=4y>0$, 故 $|f(z)|<1$. 反解 $z=\ii(1+w)/(1-w)$; 若 $|w|<1$, 则 $\Im z=(1-|w|^2)/|1-w|^2>0$, 因而双射.

(b) $z=\ee^{\ii\theta}\ne1$ 时直接计算实部为 $-1/2$; 更完整地, $w=(z+2)/(z-1)$ 反解为 $z=(w+2)/(w-1)$. 条件 $|z|=1$ 等价于 $|w+2|=|w-1|$, 即 $\Re w=-1/2$. 点 $z=1$ 是极点, 在扩充平面中对应直线的无穷远点; 不能将其当作有限值代入.
\end{solution}
\begin{exercise}\label{ass:pex1-4}
\textbf{原题 Pex1.4。} (a) 用 Cauchy--Riemann 方程证明 $\ee^z$ 解析. (b) 证明 $\bar z$ 不解析. (c) 给出 $z$ 平面中的一个“区域”, 使 $w=z^3$ 在其上一一映到整个 $w$ 平面. (d) 选择 $z^{1/3}$ 的一个解析分支及其定义区域, 使像包含在(c)所选区域内.
\end{exercise}
\begin{solution}
据官方解答翻译并校订(c)的用语. (a) $u=\ee^x\cos y,v=\ee^x\sin y$, 偏导连续且 $u_x=v_y, u_y=-v_x$, 故在全平面解析. (b) $u=x,v=-y$ 使 $u_x=1\ne-1=v_y$, 故处处不可复微分.

(c) 若“区域”指开连通集, 原要求不能成立: 满射必须含0, 但 $z^3$ 在0附近不单射. 官方实际选择的是半开扇形集合 $S=\{0\}\cup\{r\ee^{\ii\theta}:r>0,0\le\theta<2\pi/3\}$, 不是开区域. 任一非零 $w=\rho\ee^{\ii\phi}$ 取唯一 $0\le\phi<2\pi$, 在 $S$ 中的唯一原像为 $\rho^{1/3}\ee^{\ii\phi/3}$; 0的原像为0. 这给出集合意义下的双射.

(d) 在 $D=\C\setminus[0,\infty)$ 上取 $0<\arg z<2\pi$, 定义 $g(z)=\exp((\log|z|+\ii\arg z)/3)$. 所选对数解析, 故 $g$ 解析; 像为 $0<\arg w<2\pi/3$, 包含于 $S$. 下图实线边包含、虚线边不包含; 原点单独包含.
\begin{center}\begin{tikzpicture}[scale=.8]
\fill[main!10](0,0)--(3,0) arc(0:120:3)--cycle;
\draw[->](-1.8,0)--(3.4,0) node[right]{$\Re z$};\draw[->](0,-.4)--(0,3.2) node[above]{$\Im z$};
\draw[thick](0,0)--(3,0);\draw[thick,dashed](0,0)--(-1.5,2.598);\fill(0,0)circle(2pt);\node at (.8,1.4){$S$};\node at (1,.42){$2\pi/3$};
\end{tikzpicture}\end{center}
\end{solution}
\begin{exercise}\label{ass:pex1-5}
\textbf{原题 Pex1.5。} (a) 单位正方形 $C$ 上求 $\int_C x\dd z$. (b) 单位圆周上求 $\int_C |z|^{-1}\dd z$. (c) 单位圆周上求 $\int_C z\cos(z^2)\dd z$. (d) 画出 $D=\C\setminus\{x+\ii\sin x:x\ge0\}$, 判断是否单连通、能否定义对数分支. (e) 在 $|z|=3$ 上求 $\int_C z^2/(z^4-1)\dd z$. (f) 简单闭曲线 $C$ 上 $\int_C\ee^z/z^2\dd z$ 是否总为0? (g) 求 $\int_{-\infty}^{\infty}(x^4+16)^{-1}\dd x$.
\end{exercise}
\begin{solution}
据官方解答翻译并补全(d)的拓扑论证、(f)的边界条件, 校正(g)的弧估计范围. (a) 正方形顶点按 $0,1,1+\ii,\ii$ 顺序走, 四边贡献分别为 $1/2,\ii,-1/2,0$, 和为 $\ii$. (b) 圆上 $|z|=1$, 因而积分为 $\int_C\dd z=0$. (c) 被积函数整, 或取原函数 $\sin(z^2)/2$, 积分为0.

(d) 同胚 $H(x,y)=(x,y-\sin x)$ 将删去的曲线变为非负实轴, 故 $D$ 同胚于割平面, 单连通. $0\notin D$, 因而 $1/z$ 在此有原函数, 由其原函数可以构造解析对数. 不能把 $H$ 误当保角映射; 此处只用拓扑性质.
\begin{center}\begin{tikzpicture}[xscale=.8,yscale=.7,>=Stealth]
\draw[->](-1,0)--(7.1,0) node[right]{$\Re z$};\draw[->](0,-1.5)--(0,1.8) node[above]{$\Im z$};
\draw[thick,main,domain=0:6.6,samples=100] plot(\x,{sin(\x r)});\fill[main](0,0)circle(2pt);\node[above] at(4.4,1.1){删去的曲线};
\end{tikzpicture}\end{center}

(e) 四极点 $a\in\{1,-1,\ii,-\ii\}$ 的留数为 $a^2/(4a^3)=1/(4a)$, 和为0, 积分为0. (f) 若 $0\notin C$, 积分为 $2\pi\ii\,\operatorname{Ind}(C,0)$, 因为0处留数为1. 正向包围0得 $2\pi\ii$, 不包围得0; 若曲线经过0, 通常积分未定义, 原题应先排除此情形.

(g) 上半圆内两极点为 $a=2\ee^{\ii\pi/4}$、$b=2\ee^{3\ii\pi/4}$, 留数分别 $1/(4a^3),1/(4b^3)$, 和为 $-\ii\sqrt2/32$. 对 $R>2$, 弧积分绝对值不超过 $\pi R/(R^4-16)\to0$. 留数定理给结果 $2\pi\ii(-\ii\sqrt2/32)=\pi\sqrt2/16$. 官方把弧估计条件写成 $R>1$, 应改为 $R>2$.
\end{solution}
\begin{exercise}\label{ass:pex1-6}
\textbf{原题 Pex1.6。} 若 $f$ 整且处处 $|f(z)|>1$, 证明 $f$ 为常数.
\end{exercise}
\begin{solution}据官方解答翻译. $f$ 无零点, $1/f$ 整且模小于1. Liouville 定理给 $1/f$ 常数, 该常数非零, 故 $f$ 常数.\end{solution}
\begin{exercise}\label{ass:pex1-7}
\textbf{原题 Pex1.7。} $f$ 在闭圆盘 $|z-z_0|\le r$ 上解析且 $|f|$ 为常数, 证明 $f$ 在圆盘上为常数.
\end{exercise}
\begin{solution}据官方解答翻译并明确条件. “闭圆盘上解析”按其开邻域上解析理解. 若 $r>0$, 模在每个内部点取最大值, 最大模原理给开圆盘上 $f$ 常数, 连续性将结论延至边界. $r=0$ 时集合只有一点, 结论自然成立.\end{solution}
\begin{exercise}\label{ass:pex1-8}
\textbf{原题 Pex1.8。} (a) $f(z)=\ee^{\cos z}z^2$, $A=\{|z-5|\le2\}$. 证明其最大模、最小模均在边界取得, 提示考察 $1/f$. (b) 整函数 $f$ 若 $f^{(4)}$ 全平面有界, 证明 $f$ 为次数至多4的多项式. (c) $1/z^2$ 在无穷远趋于0但非常数, 是否违反 Liouville 定理?
\end{exercise}
\begin{solution}
据官方解答翻译. (a) $f$ 整, $A$ 不含0, 指数也不为零, 故 $f$ 与 $1/f$ 在 $A$ 的邻域上解析. 分别应用最大模原理, 得 $|f|$ 最大值及 $1/|f|$ 最大值在边界取得, 后者等价于 $|f|$ 最小值在边界取得. (b) Liouville 定理给 $f^{(4)}=c$. 连续积分四次, 或由 Taylor 系数 $f^{(n)}(0)=0$ ($n\ge5$), 得次数至多4. (c) 不违反, 因为 $1/z^2$ 在0不解析, 不是整函数. 本题(a)--(c)从作业4重列, 保留原考试编号以供核对.
\end{solution}
\begin{exercise}\label{ass:pex1-9}
\textbf{原题 Pex1.9。} 证明 $\int_0^\pi\ee^{\cos\theta}\cos(\sin\theta)\dd\theta=\pi$. 提示在单位圆周上积分 $\ee^z/z$.
\end{exercise}
\begin{solution}
据官方解答翻译. Cauchy 公式给 $\int_{|z|=1}\ee^z/z\dd z=2\pi\ii$. 参数化 $z=\ee^{\ii\theta}$ 得 $\ii\int_0^{2\pi}\ee^{\cos\theta+\ii\sin\theta}\dd\theta=2\pi\ii$. 取虚部得全周期的实积分为 $2\pi$. 被积函数在 $\theta\mapsto2\pi-\theta$ 下不变, 故 $[0,\pi]$ 的积分占一半.
\end{solution}
\begin{exercise}\label{ass:pex1-10}
\textbf{原题 Pex1.10。} (a) $f$ 在单连通区域 $A$ 上解析, 简单闭曲线 $\gamma\subset A$, $z_0\in A\setminus\gamma$. 用 Cauchy 公式证明 $\int_\gamma f'(z)/(z-z_0)\dd z=\int_\gamma f(z)/(z-z_0)^2\dd z$. (b) 挑战: 删除 $A$ 单连通的条件, 再证明相等.
\end{exercise}
\begin{solution}
据官方解答翻译并校正(a)漏掉的外部情形. (a) 单连通使 $\gamma$ 的内部也位于 $A$. 若 $z_0$ 在内部, 两积分均为 $2\pi\ii f'(z_0)$（正向）; 若 $z_0$ 在外部, 两被积函数在内部解析, 两积分均为0. 官方只写了前一种值.

(b) 令 $g(z)=f(z)/(z-z_0)$, 它在 $\gamma$ 的邻域上是单值解析函数, $g'=f'/(z-z_0)-f/(z-z_0)^2$. 对闭曲线积分导数为 $g(\gamma(1))-g(\gamma(0))=0$, 故所求相等. 这并不要求曲线内部都属于 $A$.
\end{solution}
\begin{exercise}\label{ass:pex1-11}
\textbf{原题 Pex1.11。} 计算以下正向圆周积分: (a) $|z|=1$ 上 $\cos z/z$; (b) $|z|=1$ 上 $\sin z/z$; (c) $|z|=2$ 上 $z^2/(z-1)$; (d) $|z|=1$ 上 $\ee^z/z^2$; (e) $|z|=2$ 上 $(z^2-1)/(z^2+1)$; (f) $|z|=2$ 上 $1/(z^2+z+1)$. 每式均对 $\dd z$ 积分.
\end{exercise}
\begin{solution}
据官方解答翻译并列出抵消步骤. (a) Cauchy 公式得 $2\pi\ii\cos0=2\pi\ii$. (b) 得 $2\pi\ii\sin0=0$, 实际0处可去. (c) 得 $2\pi\ii\cdot1^2=2\pi\ii$. (d) 导数公式得 $2\pi\ii(\ee^z)'|_0=2\pi\ii$; 这是第5题(f)的包围0特例.

(e) $\pm\ii$ 处留数为 $-2/(2\ii)$、$-2/(-2\ii)$, 和为0, 故积分0. (f) 两根 $a,b=(-1\pm\ii\sqrt3)/2$ 都在圆内, 留数分别为 $1/(a-b)$ 与 $1/(b-a)$, 和为0, 故积分0.
\end{solution}
\begin{exercise}\label{ass:pex1-12}
\textbf{原题 Pex1.12。} 整函数 $f$ 满足 $f(z)/z\to0$ ($z\to\infty$), 证明 $f$ 常数. 原卷准许使用其称作 Morera 定理的孤立点可去版本: 在 $A\setminus\{z_0\}$ 解析且在 $A$ 连续的函数在 $A$ 解析.
\end{exercise}
\begin{solution}
据官方解答翻译并补全零点定义. 令 $g(z)=(f(z)-f(0))/z$ ($z\ne0$), $g(0)=f'(0)$. $g$ 连续且0处可去, 因而整. 已知条件给 $g(z)\to0$ 于无穷远. 外部有界, 内部由紧集连续性有界, 故 $g$ 全平面有界. Liouville 给 $g$ 常数, 极限给此常数0, 所以 $f(z)=f(0)$. 此处只需可去奇点定理; 原卷给出的命题是 Morera 定理的一项推论.
\end{solution}
